% =============================================================================
% БИЛЕТ 22
% =============================================================================

\ticket{22}{}

\question{Теорема Римана о перестановке членов условно сходящегося ряда}

\begin{theorem}[Риман]
  Пусть ряд $\sum a_n$ сходится условно (т.\,е.\ $\sum a_n$ сходится, но $\sum |a_n|$ расходится).
  Тогда для любого $A \in \mathbb{R}$ существует перестановка $k: \mathbb{N} \to \mathbb{N}$ такая,
  что $\sum_{n=1}^{\infty} a_{k(n)} = A$.
  \begin{proof}
    Разделим члены ряда на положительные $p_n$ и отрицательные $q_n$
    (занумерованные в порядке их появления).

    \textbf{Оба вспомогательных ряда расходятся.}
    Если бы $\sum p_n < \infty$, то и $\sum|a_n| = \sum p_n - \sum q_n < \infty$ --- противоречие.
    Аналогично для $\sum q_n$. Значит $\sum p_n = +\infty$ и $\sum q_n = -\infty$.

    \textbf{Конструкция перестановки.}
    Будем строить новый ряд: пока его частичная сумма меньше $A$, берём элемент из $p_n$;
    как только превзошли $A$ --- берём из $q_n$; и так далее.

    Этот ряд сходится к $A$: поскольку $a_n \to 0$ (необходимое условие сходимости исходного ряда),
    «перелёты» выше и ниже $A$ стремятся к нулю, и начиная с некоторого $k_i$ мы остаёмся
    в $\varepsilon$-окрестности $A$. \qedhere
  \end{proof}
\end{theorem}

\question{Свойства функций, непрерывных на компакте: ограниченность, достижимость sup и inf, равномерная непрерывность}

\begin{theorem}[Ограниченность и достижимость экстремумов]
  Пусть $f: U \to \mathbb{R}$ непрерывна, $U \subset \mathbb{R}^n$ компактно.
  Тогда $f$ ограничена и достигает своего максимума и минимума на $U$.
  \begin{proof}
    \textbf{Ограниченность.} Допустим, $f$ не ограничена сверху: $\exists\{x_n\}$:
    $f(x_n) > n$. По теореме Больцано--Вейерштрасса $\exists$ подпоследовательность
    $x_{n_k} \to x_0 \in U$ (компактность). По непрерывности $f(x_{n_k}) \to f(x_0) \in \mathbb{R}$,
    но $f(x_{n_k}) > n_k \to \infty$ --- противоречие.

    \textbf{Достижимость $\sup$.}
    По определению $\sup$: $\forall\varepsilon>0\;\exists x\in U$: $M - \varepsilon < f(x)$.
    Строим последовательность $\{x_n\}$: $f(x_n) \in U_{\varepsilon_n}(M)$, $\varepsilon_n \to 0$.
    Выделяем сходящуюся подпоследовательность $x_{n_k} \to x_0 \in U$.
    По непрерывности $f(x_0) = \lim f(x_{n_k}) = M$. \qedhere
  \end{proof}
\end{theorem}

\begin{theorem}[Теорема Кантора об равномерной непрерывности]
  Если $f: U \to \mathbb{R}$ непрерывна на компакте $U \subset \mathbb{R}^n$,
  то $f$ равномерно непрерывна:
  \[
    \forall\varepsilon>0\;\exists\delta>0\;\forall x,y\in U:\; \rho(x,y)<\delta
    \;\Rightarrow\; |f(x)-f(y)|<\varepsilon.
  \]
  \begin{proof}
    От противного: $\exists\varepsilon_0>0$: $\forall\delta>0\;\exists x,y\in U$:
    $|f(x)-f(y)|\geq\varepsilon_0$ при $\rho(x,y) < \delta$.

    Берём $\delta = 1/n$ и получаем последовательности $\{x_n\}$, $\{y_n\}$ с $\rho(x_n,y_n) \leq 1/n$
    и $|f(x_n)-f(y_n)|\geq\varepsilon_0$.
    По компактности $\exists$ подпоследовательность $x_{n_k} \to x_0 \in U$.
    Так как $\rho(x_{n_k}, y_{n_k}) \leq 1/n_k \to 0$, то $y_{n_k} \to x_0$ тоже.
    По непрерывности: $f(x_{n_k}) \to f(x_0)$ и $f(y_{n_k}) \to f(x_0)$,
    откуда $|f(x_{n_k})-f(y_{n_k})|\to 0$ --- противоречие с $\geq\varepsilon_0$. \qedhere
  \end{proof}
\end{theorem}
